\documentclass[10pt,a4paper]{article}

% Margin lebih sempit agar space lebih banyak
\usepackage[left=1cm, right=1cm, top=1cm, bottom=1cm]{geometry}
\usepackage{titlesec}
\usepackage{enumitem}
\usepackage{hyperref}
\usepackage{tabularx}
\usepackage{xcolor}
\usepackage{colortbl}
\usepackage{array} % untuk tabular yang rapi
\usepackage{pifont} % di preamble
\usepackage{hyperref}
\usepackage[indonesian]{babel}
%---------------------------------------------------------------------------------

\usepackage{titlesec} % Allows creating custom \sections

% Format of the section titles
\titleformat{\section}
{\color[HTML]{800020}\bfseries\scshape\large\raggedright} % Biru profesional
{}{0em}{}
\titlespacing{\section}{0pt}{10pt}{4pt}
%----------------------------------------------------

% Gunakan fontspec untuk font modern seperti Raleway
\usepackage[default]{raleway}

% Hilangkan spasi antar paragraf
\setlength{\parskip}{0pt}
\setlength{\parindent}{0pt}

\begin{document}

\begin{center}
  {\huge{Rendy Ramadhana Wahyudi}}\\[0.25em]
    
  {\scshape \normalsize Instrumentation \& Field Engineering \quad | \quad Energy Service Readiness \quad | \quad Oilfield Enthusiast}\\[0.25em]

  
  \rule{\textwidth}{0.4pt}\\[0.25em]
  \normalsize
  Bandung, West Java \quad | \quad 
  \href{mailto:rendy.rw.rw@gmail.com}{rendy.rw.rw@gmail.com} \quad | \quad
  \href{https://rendy-rw.github.io/}{rendy-rw.github.io} \quad | \quad
  \href{https://www.linkedin.com/in/rendywahyudi}{linkedin.com/in/rendywahyudi}
\end{center}

\section*{Professional Summary}
Instrumentation-focused graduate with strong field experience in geophysics and advanced sensor systems. Dedicated to integrating hardware-software solutions for efficient and reliable field operations in energy services. Comfortable working in remote conditions, committed to safety, operational excellence, and dependable team coordination in dynamic field environments.

\section*{Core Competencies}
\renewcommand{\arraystretch}{1.2} % lebih padat
\begin{tabularx}{\textwidth}{>{\raggedright\arraybackslash}X >{\raggedright\arraybackslash}X >{\raggedright\arraybackslash}X}
  \rowcolor[HTML]{DCDCDC} 
  \ding{51} Instrumentation & \ding{51} Data Acquisition & \ding{51} System Maintenance \\
  \rowcolor[HTML]{FFFFFF} 
  \ding{51} Sensors & \ding{51} Control System & \ding{51} Signal Processing \\
  \rowcolor[HTML]{DCDCDC} 
  \ding{51} Electronics & \ding{51} Scientific Computing & \ding{51} Embedded Programming \\
\end{tabularx}

\section*{Professional Experience}

\setlength{\fboxsep}{0pt}
\colorbox[HTML]{DCDCDC}{%
  \parbox{\linewidth}{%
    \textbf{Geophysics Technician (Contract-Based)}, PT. Surya Brinka Persada, Bandung \quad | \quad Sep 2024 – Present
  }%
}

Led 2D resistivity geophysical surveys to support groundwater potential assessments across West Java.

\begin{itemize}[left=1.5em, noitemsep, topsep=0pt]
  \item Coordinated a three-member field team to conduct geophysical surveys at 16 dispersed locations, optimizing technical role assignments to significantly reduce survey duration and operational costs.
  \item Completed all projects ahead of schedule, delivering acquisition data with less than 10\% error—accepted by stakeholders without the need for revision.
\end{itemize}


\setlength{\fboxsep}{0pt}
\colorbox[HTML]{DCDCDC}{%
  \parbox{\linewidth}{%
    \textbf{Instrumentation Engineer (Contract-Based)}, Geophysics Research Team, UNPAD, Sumedang \quad | \quad Jul – Aug 2023
  }%
}

Developed systems under tight deadlines, applying various approaches to enhance operational effectiveness.

\begin{itemize}[left=1.5em, noitemsep, topsep=0pt]
  \item Initiated Instrumentation performance optimization using open-source numerical simulation software, resulting in up to 50\% savings in cost and material use.
  \item Designed and built a functional prototype that served as the foundation for measurable and sustainable further research development.
\end{itemize}

\section*{Leadership \& Field Coordination Experience}

\noindent
\setlength{\fboxsep}{0pt}%
\colorbox[HTML]{DCDCDC}{%
  \parbox{\linewidth}{%
    \textbf{Village Empowerment Program Coordinator}, HIMA Geofisika PEDRA, UNPAD \quad | \quad Sep 2021 – Feb 2022%
  }%
}

Coordinated a 10-member team for a long-term community development program reaching over 200 residents.

\begin{itemize}[left=1.5em, noitemsep, topsep=0pt]
  \item Organized earthquake preparedness training using the CBDRR approach and conducted emergency response mock drills to improve information dissemination effectiveness.
  \item Successfully increased community disaster awareness by 80\% based on post-activity surveys.
\end{itemize}

\section*{Field Readiness \& Safety}

\begin{itemize}[left=1.5em, noitemsep, topsep=0pt]
  \item Eager to learn and gain hands-on experience in remote, challenging environments.
  \item Experienced working independently under minimal supervision, with a strong focus on safety standards.
  \item Physically and mentally prepared for long hours and demanding field conditions.
\end{itemize}

\section*{Education}

\setlength{\fboxsep}{0pt}%
\colorbox[HTML]{DCDCDC}{%
  \parbox{\linewidth}{%
    \textbf{Master of Science in Physics}, Universitas Padjadjaran, West Java \quad | \quad Jan 2023 – Aug 2025 (Expected)%
  }%
}
\begin{itemize}[left=1.5em, noitemsep, topsep=0pt]
  \item \textbf{Capstone Project}: \textit{Development of a disaster monitoring system integrating environmental sensors, real-time data acquisition, and embedded control for regional early-warning applications.}
\end{itemize}

\setlength{\fboxsep}{0pt}%
\colorbox[HTML]{DCDCDC}{%
  \parbox{\linewidth}{%
    \textbf{ Bachelor of Science in Geophysical}, Universitas Padjadjaran, West Java \quad | \quad Aug 2019 – Jul 2023%
  }%
}
{\raggedright\sloppy
\begin{itemize}[left=1.5em, noitemsep, topsep=0pt]
  \item \textbf{Capstone Project}: \textit{Prototype development of IoT-based distributed seismic sensor network}, published at \linebreak DOI: \href{https://doi.org/10.2991/978-2-38476-283-5_7}{10.2991/978-2-38476-283-5\_7}.
\end{itemize}
}

\section*{Certificates \& Grants}

\textbf{Graduate Scholarship Awardee}, JFLS, Disdik Jabar \quad | \quad Feb 2025 \\
\textbf{Graduate Research Grant}, Kemdiktisaintek \quad | \quad Aug 2024 \\
\textbf{TOEFL ITP Score: 527, Conversational English}, LBI FIB Universitas Indonesia \quad | \quad Jul 2024

\section*{Technical Skills}

\textbf{Instrumentation:} Sensor Integration, STM32, ESP32 \quad | \quad
\textbf{Embedded Systems:} C++, Arduino, STM32CubeIDE, KiCad   \quad | \quad
\textbf{Signal \& Data Processing:} Python (NumPy, SciPy), MATLAB, Machine learning, Big Data \quad | \quad
\textbf{Simulation \& Testing Tools:} Proteus, LabVIEW \quad | \quad
\textbf{Documentation:} LaTeX, MS Word \quad | \quad
\textbf{Field Tools \& OS:} Linux (basic terminal ops), Windows, Microosft Teams, Primavera P6.
\end{document}
